\documentclass[12pt]{article}

% ==================== Core math ====================
\usepackage{amsmath, amssymb, amsthm}

% ==================== Tables ====================
\usepackage{booktabs}
\usepackage{array}
\usepackage{longtable}
\usepackage{aliascnt}

% ==================== Diagrams ====================
\usepackage{tikz-cd}
\usepackage{tikz}
\usetikzlibrary{arrows.meta, positioning, calc, fit, backgrounds}

% ==================== References ====================
\usepackage{hyperref}
\usepackage{cleveref}

% ==================== Graphics ====================
\usepackage{graphicx}

% ==================== Page layout ====================
\usepackage[margin=1in]{geometry}
% Allow TeX a little extra stretch to eliminate residual sub-1em overfull boxes.
\setlength{\emergencystretch}{1.5em}

% ==================== GrokRxiv sidebar ====================
% (Modern LaTeX kernel shipout hooks are used below instead of the deprecated
%  everypage package.)
\usepackage{xcolor}

\hypersetup{
  colorlinks=true,
  linkcolor=blue!55!black,
  citecolor=green!45!black,
  urlcolor=blue!60!black,
  pdftitle={Algebraic Topology: Conserved Global Information},
  pdfauthor={Matthew Long}
}

% ==================== Theorem environments ====================
% All environments share the theorem counter for unified numbering, but via
% aliascnt so that cleveref labels each one correctly (Definition/Proposition/...
% rather than the counter name "Theorem").
\newtheorem{theorem}{Theorem}[section]

\newaliascnt{proposition}{theorem}
\newtheorem{proposition}[proposition]{Proposition}
\aliascntresetthe{proposition}

\newaliascnt{lemma}{theorem}
\newtheorem{lemma}[lemma]{Lemma}
\aliascntresetthe{lemma}

\newaliascnt{corollary}{theorem}
\newtheorem{corollary}[corollary]{Corollary}
\aliascntresetthe{corollary}

\theoremstyle{definition}
\newaliascnt{definition}{theorem}
\newtheorem{definition}[definition]{Definition}
\aliascntresetthe{definition}

\newaliascnt{example}{theorem}
\newtheorem{example}[example]{Example}
\aliascntresetthe{example}

\theoremstyle{remark}
\newaliascnt{remark}{theorem}
\newtheorem{remark}[remark]{Remark}
\aliascntresetthe{remark}

% Explicit cleveref names (belt-and-suspenders with aliascnt).
\crefname{theorem}{Theorem}{Theorems}
\Crefname{theorem}{Theorem}{Theorems}
\crefname{proposition}{Proposition}{Propositions}
\Crefname{proposition}{Proposition}{Propositions}
\crefname{lemma}{Lemma}{Lemmas}
\Crefname{lemma}{Lemma}{Lemmas}
\crefname{corollary}{Corollary}{Corollaries}
\Crefname{corollary}{Corollary}{Corollaries}
\crefname{definition}{Definition}{Definitions}
\Crefname{definition}{Definition}{Definitions}
\crefname{example}{Example}{Examples}
\Crefname{example}{Example}{Examples}
\crefname{remark}{Remark}{Remarks}
\Crefname{remark}{Remark}{Remarks}

% ==================== Macros ====================
\newcommand{\Bord}{\mathrm{Bord}}
\newcommand{\Vect}{\mathrm{Vect}}
\newcommand{\Top}{\mathrm{Top}}
\newcommand{\hTop}{\mathrm{hTop}}
\newcommand{\GrAb}{\mathrm{GrAb}}
\newcommand{\Ab}{\mathrm{Ab}}
\newcommand{\Man}{\mathrm{Man}}
\newcommand{\Dom}{\mathsf{Dom}}
\newcommand{\Phys}{\mathsf{Phys}}
\newcommand{\Math}{\mathsf{Math}}
\newcommand{\Rep}{\mathsf{Rep}}
\newcommand{\Real}{\mathrm{Real}}
\newcommand{\Obs}{\mathrm{Obs}}
\newcommand{\Info}{\mathrm{Info}}
\newcommand{\ind}{\mathrm{ind}}
\newcommand{\ch}{\mathrm{ch}}
\newcommand{\Ahat}{\widehat{A}}
\newcommand{\id}{\mathrm{id}}
\newcommand{\im}{\operatorname{im}}
\newcommand{\Hom}{\operatorname{Hom}}
\newcommand{\Aut}{\operatorname{Aut}}
\newcommand{\Stab}{\operatorname{Stab}}
\newcommand{\Zm}{\mathbb{Z}}
\newcommand{\Rm}{\mathbb{R}}
\newcommand{\Cm}{\mathbb{C}}
\newcommand{\Qm}{\mathbb{Q}}
\newcommand{\Umone}{\mathrm{U}(1)}
\newcommand{\sfd}{^{\mathrm{fd}}}
\DeclareMathOperator{\rk}{rk}
\newcommand{\Sfont}{\textnormal{\sffamily\bfseries S}}
\newcommand{\Hfont}{\textnormal{\sffamily\bfseries H}}
\newcommand{\Pfont}{\textnormal{\sffamily\bfseries P}}
\newcommand{\st}[1]{\ensuremath{[\,#1\,]}}

% GrokRxiv rotated sidebar (page 1 only), via the modern shipout hook.
\definecolor{grokgray}{RGB}{110,110,110}
\AddToHook{shipout/foreground}{%
  \ifnum\value{page}=1
    \begin{tikzpicture}[remember picture, overlay]
      \node[
        rotate=90,
        anchor=south,
        font=\Large\sffamily\bfseries\color{grokgray},
        inner sep=0pt
      ] at ([xshift=38pt, yshift=0.52\paperheight]current page.south west)
      {GrokRxiv:2026.07.algebraic-topology-conserved-information\quad
       [\,math.AT\,]\quad
       04 Jul 2026};
    \end{tikzpicture}
  \fi
}

\title{\textbf{Algebraic Topology: Conserved Global Information}\\[0.4em]
\large Part IV of \emph{A Math$\to$Physics Representation Library}}

\author{Matthew Long \\
\textit{The YonedaAI Collaboration} \\
\textit{YonedaAI Research Collective} \\
Chicago, IL \\
\texttt{research@yonedaai.com} $\cdot$ \url{https://yonedaai.com}}

\date{04 July 2026}

\begin{document}
\maketitle

\begin{abstract}
This is Part~IV of a seven-part \emph{modular} program that formalizes distinct
mathematical domains as physical-representation modules within a single
representation-stack framework. Building on Part~I (the representation stack and
realization pipeline $\Obs_\alpha(M)=\Obs(\Real_\alpha(\Phi(M)))$), Part~II
(motives, periods, amplitudes), and Part~III (algebraic geometry, spaces of
physical possibility), we develop the thesis that \emph{algebraic topology is the
grammar of conserved global information}. We treat (co)homology as the module of
observables that are stable under continuous deformation, characteristic classes
as quantized topological information carried by gauge bundles, and topological
quantum field theory (TQFT) as the first fully rigorous, status-$\Sfont$
realization of the Part~I pipeline in \emph{functorial} form over a category of
cobordisms. Our main results are: (T1) that an $n$-dimensional TQFT
$Z\colon\Bord_n\to\Vect_k$ is exactly a representation entry in the sense of
Part~I, with its symmetric-monoidal structure discharging the Decomposition
axiom; (T2) a statement of the cobordism hypothesis as a classification theorem
for topological realization channels; (T3) a reading of the Atiyah--Singer index
theorem $\ind(D)=\int_X \Ahat(TX)\wedge\ch(E)$ as a realization-pipeline identity
$\text{observable}=\Real(\text{abstract structure})$; (T4) that Dijkgraaf--Witten
theory instantiates the Locality/descent axiom through group cohomology; and (T5)
a bordism-theoretic classification of anomalies and symmetry-protected
topological phases. We prove homotopy invariance and functoriality of homology
in full, formalize a \emph{conserved global information functor}, and carry the
$\Sfont/\Hfont/\Pfont$ epistemic status discipline of the parent program
throughout. Companion Haskell code (compiling under GHC) and a best-effort Lean
sketch encode chain complexes, homology, characteristic classes, and the TQFT
axioms. Part~IV is the point at which category theory first appears as the
ambient grammar; it thereby seeds Part~V.
\end{abstract}

\tableofcontents

% =====================================================================
\section{Introduction}
\label{sec:intro}

\subsection{The modular program and where Part~IV sits}

The parent program, \emph{A Math$\to$Physics Representation Library}, advances a
single organizing claim: many physical quantities are not merely \emph{described
by} mathematics but are \emph{obtained as realizations} of structured
mathematical information. The organizing device is the \emph{realization
pipeline}
\begin{equation}
\label{eq:pipeline}
\begin{gathered}
  M \;\xrightarrow{\;\Phi\;}\; \Phi(M)
    \;\xrightarrow{\;\Real_\alpha\;}\; \text{physical representation}
    \;\xrightarrow{\;\Obs\;}\; \text{measurement data},\\[0.4em]
  \Obs_\alpha(M) = \Obs\bigl(\Real_\alpha(\Phi(M))\bigr),
\end{gathered}
\end{equation}
together with a three-level epistemic status system that every module must carry:
\begin{center}
\small
\begin{tabular}{c l p{5.1cm}}
\toprule
Label & Meaning & Canonical example\\
\midrule
$\Sfont$ & standard use in mathematics or physics & TQFT as a functor $\Bord_n\to\Vect$\\
$\Hfont$ & strong heuristic dictionary entry & motive as ``functional essence''\\
$\Pfont$ & speculative philosophical ontology & reality as motivic-categorical realization\\
\bottomrule
\end{tabular}
\end{center}
Composite labels ($\Sfont/\Hfont$, $\Hfont/\Pfont$) appear when a construction is
standard as pure mathematics but heuristic or speculative in its \emph{physical}
reading. We use the labels verbatim; they are part of the formalism, not
decoration.

The program is explicitly \emph{modular} rather than unified: each part is
readable and defensible on its own, and composition is exhibited as a hierarchical
building process. The ladder proceeds
{\small
\[
  \underbrace{\text{I}}_{\text{stack}}\to
  \underbrace{\text{II}}_{\text{periods}}\to
  \underbrace{\text{III}}_{\text{geometry}}\to
  \underbrace{\text{IV}^{\!\star}}_{\text{topology}}\to
  \underbrace{\text{V}}_{\text{cat./HoTT}}\to
  \underbrace{\text{VI}}_{\text{stacks/gauge/QHA}}\to
  \text{Synth.},
\]}%
where $\text{IV}^{\star}$ denotes the present paper.
Part~III equips the ``spaces of physical possibility'' (varieties, moduli) with a
geometry. Part~IV equips those same spaces with \emph{topological invariants}
stable under continuous deformation, and it is the first place where a
\emph{genuinely functorial, status-$\Sfont$} instance of the pipeline
\eqref{eq:pipeline} appears: a topological quantum field theory is literally a
functor. This is the historical and conceptual bridge to Part~V, where
functoriality is abstracted into the universal grammar of categories and homotopy
type theory.

\subsection{Thesis of Part~IV}

We defend, formalize, and where possible prove the slogan
\begin{center}
\emph{algebraic topology is the grammar of conserved global information.}
\end{center}
Concretely: a homology class is a conserved extended structure; a cohomology
class is a stable observable insensitive to local deformation; characteristic
classes encode quantized topological information carried by bundles; bordism
groups and generalized cohomology classify topological phases and anomalies; and
TQFT packages all of this as a symmetric monoidal functor from cobordisms to an
algebraic target. The recurring physical motif is \emph{conservation}: what
survives continuous deformation is exactly what algebraic topology measures, and
this is the topological face of the parent program's Decomposition and Locality
axioms.

\subsection{Contributions}

\begin{enumerate}
  \item \textbf{A representation-stack account of (co)homology}
    (\Cref{sec:homology}). We define a \emph{conserved global information functor}
    $H_\bullet\colon\hTop\to\GrAb$ and prove homotopy invariance and
    functoriality in full, giving a precise sense in which topological observables
    are \emph{conserved} under physical deformation (\Cref{thm:homotopy-inv,thm:cons-functor}).
  \item \textbf{TQFT as a status-$\Sfont$ realization entry} (\Cref{sec:tqft}).
    We show an $n$-dimensional TQFT is exactly a Part~I representation entry, and
    that symmetric monoidality is the Decomposition axiom over cobordisms
    (\Cref{thm:tqft-entry}); we state the cobordism hypothesis as a classification
    theorem for topological realization channels (\Cref{thm:cobordism-hyp}).
  \item \textbf{The index theorem as a pipeline identity} (\Cref{sec:char}).
    We reframe Atiyah--Singer as an instance of
    $\text{observable}=\Real(\text{abstract structure})$
    (\Cref{thm:index}), the strongest available status-$\Sfont$ anchor.
  \item \textbf{Discrete gauge theory and locality} (\Cref{sec:dw}). We show
    Dijkgraaf--Witten theory instantiates the Locality/descent axiom via a group
    cocycle, with triangulation independence coming from Pachner-move invariance
    (\Cref{thm:dw}).
  \item \textbf{Bordism classification of anomalies and phases}
    (\Cref{sec:bordism}). We record the Freed--Hopkins classification of
    invertible field theories and its physical reading as a status-$\Sfont$ /
    $\Sfont\!/\Hfont$ statement (\Cref{thm:spt}).
  \item \textbf{Formalization} (\Cref{sec:code}). Companion Haskell code
    (compiling under GHC) and a Lean sketch encode chain complexes, homology,
    characteristic classes and TQFT axioms, extending the parent program's
    formal-verification bridge.
\end{enumerate}

\subsection{Epistemic status of Part~IV}

Part~IV is unusually rich in status-$\Sfont$ content: TQFT, the index theorem,
Dijkgraaf--Witten theory, and the bordism classification of invertible field
theories are all standard mathematics or mathematical physics. The heuristic
($\Hfont$) content lies in the \emph{physical readings} of dictionary entries
(e.g.\ ``coboundary $=$ pure gauge''), and the speculative ($\Pfont$) content is
confined to the ontological framing inherited from the parent program. We flag
labels at every dictionary entry (\Cref{sec:dict}) and at each theorem.

\subsection{Outline}

\Cref{sec:framework} recalls the representation stack and pipeline from Part~I in
the precise form used here. \Cref{sec:homology} develops (co)homology as conserved
global information. \Cref{sec:char} treats characteristic classes and the index
theorem. \Cref{sec:tqft} develops TQFT and the cobordism hypothesis.
\Cref{sec:dw} treats Dijkgraaf--Witten theory; \Cref{sec:bordism} treats bordism
and SPT phases. \Cref{sec:results} collects the main results. \Cref{sec:code}
describes the companion code. \Cref{sec:discussion} discusses limitations and
connections to Parts~III and~V, and \Cref{sec:conclusion} concludes.
\Cref{sec:dict} reproduces the algebraic-topology dictionary with status labels.

% =====================================================================
\section{Mathematical Framework: the representation stack, recalled}
\label{sec:framework}

We recall from Part~I only what Part~IV uses, in a self-contained form.

\begin{definition}[Domains and targets]
\label{def:dom}
Let $\Dom$ be a (small) category of \emph{mathematical domains}. For each
$d\in\Dom$ let $\Math_d$ be a category (or higher category) of mathematical
structures in that domain. Let $\Phys$ be a category of \emph{physical
representation targets}: state spaces, observable algebras, process categories,
gauge moduli, quantum code spaces, and measurement outcomes. In Part~IV the
relevant domains are $\Top$ (topological spaces), $\Man_n$ (closed manifolds and
cobordisms), and the associated $\Bord_n$.
\end{definition}

\begin{definition}[Representation entry]
\label{def:entry}
A \emph{representation entry} is a quadruple $E=(M,P,\tau,\sigma)$ where
$M\in\Math_d$ is a mathematical structure, $P\in\Phys$ a physical representation,
$\tau$ a semantic translation datum (a functor, a natural transformation, or a
weaker relation), and $\sigma$ is the epistemic status label ($\Sfont$, $\Hfont$,
or $\Pfont$).
\end{definition}

\begin{definition}[Realization pipeline]
\label{def:pipeline}
A \emph{realization pipeline} on categories $M\xrightarrow{\Phi}I
\xrightarrow{\Real_\alpha}R\xrightarrow{\Obs}O$ is a composable triple of
functors; the associated \emph{observable functor} is
$\Obs_\alpha := \Obs\circ\Real_\alpha\circ\Phi$, and we write
$\Obs_\alpha(M)=\Obs(\Real_\alpha(\Phi(M)))$ as in \eqref{eq:pipeline}. The
index $\alpha$ ranges over \emph{realization channels}; in Part~IV the channels
of interest are \emph{topological}: singular/de~Rham (co)homology, characteristic
classes, TQFT evaluation, and the analytic index.
\end{definition}

Part~I isolates four axioms. We restate the two that Part~IV discharges
concretely.

\begin{definition}[Locality and descent axiom]
\label{ax:locality}
Physical representations must be locally assignable and globally coherent; a
failure of descent is an obstruction, an anomaly, or a missing boundary datum.
\end{definition}

\begin{definition}[Decomposition axiom]
\label{ax:decomp}
Compositional or period-like observables admit a decomposition operation
(boundary, coproduct, coaction, spectral sequence, factorization, or gluing)
revealing lower-level informational pieces.
\end{definition}

The master diagram of Part~I, in the form we use, is
\[
\begin{tikzcd}[column sep=1.6em, row sep=large]
M' \arrow[d, dashed, "\sim"', "\text{gauge / equiv.\ / duality}"] & & &\\
M \arrow[r, "\Phi"] & \Phi(M) \arrow[r, "\Real_\alpha"] & \text{physical rep.} \arrow[r, "\Obs"] & \text{observable}
\end{tikzcd}
\]
The dashed identification $M\sim M'$ encodes the Equivalence axiom: mathematical
descriptions related by a gauge or equivalence relation determine the same
physical content at the chosen observational level. Part~IV realizes the dashed
arrow twice: as homotopy of spaces (\Cref{sec:homology}) and as coboundary shift
of a Dijkgraaf--Witten cocycle (\Cref{sec:dw}).

\begin{remark}[Status of the framework]
\label{rem:framework-status}
\Cref{def:dom,def:entry,def:pipeline} are status-$\Sfont$ as pure category theory;
their \emph{physical} interpretation carries the label of the individual entry.
Part~IV's contribution is to populate the framework with entries whose $\tau$ is
an \emph{actual functor}, i.e.\ status-$\Sfont$ entries, in contrast to the
heuristic entries that dominate Parts~II--III.
\end{remark}

% =====================================================================
\section{(Co)homology as conserved global information}
\label{sec:homology}

We now develop the core of Part~IV: algebraic topology as the module of
observables conserved under continuous deformation. We proceed from chain
complexes to homology, prove functoriality and homotopy invariance, dualize to
cohomology, and read the results physically.

\subsection{Chain complexes and the boundary operator}

\begin{definition}[Chain complex]
\label{def:chain}
A \emph{chain complex} of abelian groups is a sequence
\[
\begin{tikzcd}[column sep=2.1em]
\cdots \arrow[r, "\partial_{n+1}"] & C_n \arrow[r, "\partial_n"] & C_{n-1}
  \arrow[r, "\partial_{n-1}"] & \cdots \arrow[r, "\partial_1"] & C_0 \arrow[r, "\partial_0"] & 0
\end{tikzcd}
\]
of abelian groups $C_n$ ($n\ge 0$) and homomorphisms $\partial_n\colon C_n\to
C_{n-1}$ satisfying the \emph{fundamental relation}
\begin{equation}
\label{eq:d2}
  \partial_{n}\circ\partial_{n+1}=0\qquad\text{for all }n .
\end{equation}
Elements of $Z_n:=\ker\partial_n$ are \emph{$n$-cycles}; elements of
$B_n:=\im\partial_{n+1}$ are \emph{$n$-boundaries}. A morphism of chain complexes
(a \emph{chain map}) $f_\bullet\colon C_\bullet\to D_\bullet$ is a family
$f_n\colon C_n\to D_n$ with $\partial^D_n f_n = f_{n-1}\partial^C_n$.
\end{definition}

\begin{lemma}[Boundaries are cycles]
\label{lem:BsubZ}
For every chain complex, $B_n\subseteq Z_n$, so $H_n:=Z_n/B_n$ is a well-defined
abelian group.
\end{lemma}

\begin{proof}
Let $b\in B_n$, say $b=\partial_{n+1}c$ for some $c\in C_{n+1}$. Then by
\eqref{eq:d2},
$\partial_n b = \partial_n\partial_{n+1}c = 0$, so $b\in\ker\partial_n=Z_n$.
Hence $B_n\subseteq Z_n$. As $C_n$ is abelian, $B_n$ is a (normal) subgroup of
$Z_n$ and the quotient $Z_n/B_n$ is an abelian group.
\end{proof}

\begin{definition}[Homology]
\label{def:homology}
The \emph{$n$-th homology group} of $C_\bullet$ is $H_n(C_\bullet):=Z_n/B_n$. The
class of a cycle $z\in Z_n$ is written $[z]$. A chain complex is \emph{exact} at
$C_n$ if $H_n(C_\bullet)=0$, i.e.\ $Z_n=B_n$.
\end{definition}

For a topological space $X$ we take $C_\bullet=C_\bullet(X)$ the singular chain
complex, with $C_n(X)$ the free abelian group on continuous maps
$\sigma\colon\Delta^n\to X$ and $\partial_n\sigma=\sum_{i=0}^n(-1)^i\,\sigma\!\mid_{[v_0,\dots,\hat v_i,\dots,v_n]}$
the alternating face sum. Then $H_n(X):=H_n(C_\bullet(X))$ is singular homology.

\begin{lemma}[Singular boundary squares to zero]
\label{lem:sing-d2}
The singular boundary satisfies $\partial_n\partial_{n+1}=0$, so $C_\bullet(X)$ is
a chain complex in the sense of \Cref{def:chain}.
\end{lemma}

\begin{proof}
It suffices to check on a single simplex $\sigma\colon\Delta^{n+1}\to X$. Writing
$\sigma_{[\hat i]}$ for the $i$-th face,
\[
  \partial_n\partial_{n+1}\sigma
  = \sum_{i=0}^{n+1}(-1)^i\,\partial_n\sigma_{[\hat i]}
  = \sum_{i=0}^{n+1}\sum_{j=0}^{n}(-1)^i(-1)^j\,(\sigma_{[\hat i]})_{[\hat j]} .
\]
The face maps satisfy the simplicial identity $(\sigma_{[\hat i]})_{[\hat j]}
=(\sigma_{[\hat j]})_{[\widehat{i-1}]}$ for $j<i$. Splitting the double sum into
$j<i$ and $j\ge i$ and applying this identity shows the two pieces cancel term by
term after reindexing, giving $0$.
\end{proof}

\subsection{Functoriality: the conserved-information functor}

\begin{proposition}[Homology is functorial]
\label{prop:H-functor}
A continuous map $f\colon X\to Y$ induces a chain map $f_\#\colon C_\bullet(X)\to
C_\bullet(Y)$ by $f_\#(\sigma)=f\circ\sigma$, hence homomorphisms
$f_*\colon H_n(X)\to H_n(Y)$ for all $n$. This assignment satisfies
$(\id_X)_*=\id$ and $(g\circ f)_*=g_*\circ f_*$; thus $H_n(-)$ is a functor
$\Top\to\Ab$.
\end{proposition}

\begin{proof}
For a singular simplex $\sigma\colon\Delta^n\to X$ and $0\le i\le n$ we have
$(f\circ\sigma)_{[\hat i]}=f\circ\sigma_{[\hat i]}$, so $f_\#$ commutes with the
face maps and hence with $\partial$; thus $f_\#$ is a chain map. A chain map sends
cycles to cycles and boundaries to boundaries, so it descends to
$f_*\colon H_n(X)\to H_n(Y)$. Functoriality of $(-)_\#$ is immediate from
$(g\circ f)\circ\sigma=g\circ(f\circ\sigma)$ and passes to homology.
\end{proof}

The physical content of Part~IV rests on the fact that $H_n$ is not merely
functorial but \emph{homotopy invariant}: it cannot distinguish continuously
deformable maps. This is precisely ``conservation under deformation.''

\begin{definition}[Chain homotopy]
\label{def:chain-homotopy}
Two chain maps $f_\bullet,g_\bullet\colon C_\bullet\to D_\bullet$ are
\emph{chain homotopic} if there is a family $P_n\colon C_n\to D_{n+1}$ with
\begin{equation}
\label{eq:chain-htpy}
  \partial^D_{n+1}P_n + P_{n-1}\partial^C_n = g_n - f_n\qquad\text{for all }n .
\end{equation}
\end{definition}

\begin{lemma}[Chain homotopy $\Rightarrow$ equality on homology]
\label{lem:chainhtpy-homology}
If $f_\bullet$ and $g_\bullet$ are chain homotopic, then $f_*=g_*$ on homology.
\end{lemma}

\begin{proof}
Let $z\in Z_n$, so $\partial z=0$. Applying \eqref{eq:chain-htpy},
\[
  (g_n-f_n)(z) = \partial_{n+1}P_n(z) + P_{n-1}\partial_n(z)
  = \partial_{n+1}\bigl(P_n(z)\bigr) + P_{n-1}(0)
  = \partial_{n+1}\bigl(P_n(z)\bigr)\in B_n .
\]
Hence $g_n(z)$ and $f_n(z)$ differ by a boundary, so $[g_n(z)]=[f_n(z)]$ in
$H_n(D_\bullet)$; that is, $g_*=f_*$.
\end{proof}

\begin{theorem}[Homotopy invariance of homology]
\label{thm:homotopy-inv}
If $f,g\colon X\to Y$ are homotopic continuous maps, then $f_*=g_*\colon
H_n(X)\to H_n(Y)$ for all $n$. Consequently a homotopy equivalence induces
isomorphisms on all homology groups, and $H_n$ descends to a functor
\[
  H_n\colon \hTop \longrightarrow \Ab
\]
on the homotopy category $\hTop$ (spaces and homotopy classes of maps).
\end{theorem}

\begin{proof}
A homotopy $F\colon X\times[0,1]\to Y$ from $f$ to $g$ induces a chain homotopy
between $f_\#$ and $g_\#$ via the \emph{prism operator}
$P_n\colon C_n(X)\to C_{n+1}(Y)$ obtained by subdividing the prism
$\Delta^n\times[0,1]$ into $(n+1)$-simplices $[v_0,\dots,v_i,w_i,\dots,w_n]$
(with $v_j=(x,0)$, $w_j=(x,1)$ vertices) and setting
$P_n(\sigma)=\sum_{i=0}^n(-1)^i F\circ(\sigma\times\id)\!\mid_{[v_0,\dots,v_i,w_i,\dots,w_n]}$.
A direct computation of $\partial P + P\partial$ using the simplicial identities
yields $g_\# - f_\#$, i.e.\ \eqref{eq:chain-htpy}. By
\Cref{lem:chainhtpy-homology}, $f_*=g_*$. If $f$ is a homotopy equivalence with
inverse $h$, then $f_*h_*=(fh)_*=(\id)_*=\id$ and similarly $h_*f_*=\id$, so
$f_*$ is an isomorphism. Since $H_n$ agrees on homotopic maps, it factors through
$\hTop$.
\end{proof}

We package this as the central object of Part~IV.

\begin{definition}[Conserved global information functor]
\label{def:cgi}
The \emph{conserved global information functor} is
\[
  H_\bullet \colon \hTop \longrightarrow \GrAb,
  \qquad X\longmapsto \bigl(H_n(X)\bigr)_{n\ge 0},
\]
valued in graded abelian groups, assembled from \Cref{thm:homotopy-inv}. A
\emph{topological observable} of $X$ is an element of $H^n(X)$ (see
\Cref{def:cohomology}); its value on a cycle is a \emph{conserved charge}.
\end{definition}

\begin{theorem}[Conservation principle]
\label{thm:cons-functor}
Let $\{X_t\}_{t\in[0,1]}$ be a continuous family of spaces realized as the fibers
of a fibration $E\to[0,1]$ over the contractible base $[0,1]$. Then for all $n$
the groups $H_n(X_t)$ are canonically isomorphic, and any topological observable
$\omega\in H^n(X_0)$ has a unique continuation $\omega_t\in H^n(X_t)$ with
$\langle\omega_t,[z_t]\rangle$ independent of $t$ for a continuously varying cycle
class $[z_t]$. In this precise sense topological observables are conserved under
continuous physical deformation.
\end{theorem}

\begin{proof}
A fibration over a contractible base is fiber-homotopy equivalent to the trivial
fibration $X_0\times[0,1]\to[0,1]$; in particular any two fibers $X_s,X_t$ are
homotopy equivalent, and parallel transport along $[0,1]$ furnishes a canonical
homotopy class of equivalence $X_0\simeq X_t$. By \Cref{thm:homotopy-inv} the
induced maps $H_n(X_0)\xrightarrow{\ \cong\ }H_n(X_t)$ are isomorphisms;
dualizing (\Cref{prop:coho-functor}) gives the corresponding isomorphisms on
$H^n$ and hence the continuation $\omega_t$. The pairing
$\langle\omega_t,[z_t]\rangle$ is preserved because both $\omega_t$ and $[z_t]$
are images of $\omega_0,[z_0]$ under the same transport isomorphism, which is
compatible with the (co)homology pairing (\Cref{def:pairing}). Thus the pairing is
constant in $t$.
\end{proof}

\subsection{Cohomology, closed and exact forms, and gauge}

\begin{definition}[Singular cohomology]
\label{def:cohomology}
For an abelian group $A$, the \emph{singular cochain complex} of $X$ is
$C^n(X;A):=\Hom(C_n(X),A)$ with coboundary
$\delta^n=\Hom(\partial_{n+1},A)$, i.e.\ $(\delta\varphi)(c)=\varphi(\partial c)$.
Then $\delta^{n+1}\delta^n=0$ (dual to \eqref{eq:d2}), and the \emph{$n$-th
cohomology group} is $H^n(X;A):=\ker\delta^n/\im\delta^{n-1}$. Cocycles are
$Z^n$, coboundaries are $B^n$.
\end{definition}

\begin{proposition}[Cohomology is a contravariant functor]
\label{prop:coho-functor}
$H^n(-;A)\colon\hTop^{\mathrm{op}}\to\Ab$ is a functor: a map $f\colon X\to Y$
induces $f^*=\Hom(f_\#,A)\colon H^n(Y;A)\to H^n(X;A)$, with $(\id)^*=\id$,
$(gf)^*=f^*g^*$, and $f^*$ depending only on the homotopy class of $f$.
\end{proposition}

\begin{proof}
$\Hom(-,A)$ is a contravariant additive functor, so it sends the chain map $f_\#$
to a cochain map $f^*$ and reverses composition; homotopy invariance is dual to
\Cref{thm:homotopy-inv} because $\Hom(-,A)$ sends a chain homotopy \eqref{eq:chain-htpy}
to a cochain homotopy.
\end{proof}

\begin{definition}[Evaluation pairing]
\label{def:pairing}
The \emph{Kronecker pairing} $\langle-,-\rangle\colon H^n(X;A)\times H_n(X)\to A$
is $\langle[\varphi],[z]\rangle:=\varphi(z)$. It is well defined: if
$\varphi=\delta\psi$ then $\varphi(z)=\psi(\partial z)=0$ for a cycle $z$, and if
$z=\partial w$ then $\varphi(z)=(\delta\varphi)(w)=0$ for a cocycle $\varphi$.
\end{definition}

The de~Rham incarnation makes the ``conserved field'' reading transparent.

\begin{definition}[de~Rham cohomology]
\label{def:derham}
For a smooth manifold $X$, let $\Omega^n(X)$ be the smooth $n$-forms with exterior
derivative $d\colon\Omega^n\to\Omega^{n+1}$, $d^2=0$. A form is \emph{closed} if
$d\omega=0$ and \emph{exact} if $\omega=d\alpha$. The \emph{de~Rham cohomology} is
$H^n_{\mathrm{dR}}(X):=\ker(d\!\mid_{\Omega^n})/\im(d\!\mid_{\Omega^{n-1}})$.
\end{definition}

\begin{proposition}[Conservation and gauge, de~Rham form]
\label{prop:derham-gauge}
Let $\omega\in\Omega^n(X)$ be a field strength. Then:
\begin{enumerate}
  \item[(i)] $d\omega=0$ (a conservation law / Bianchi identity) is the statement
    that $\omega$ is a cocycle; its class $[\omega]\in H^n_{\mathrm{dR}}(X)$ is a
    deformation-invariant observable.
  \item[(ii)] a shift $\omega\mapsto\omega+d\alpha$ leaves the class $[\omega]$
    unchanged, so $[\omega]$ is the gauge-invariant flux and an exact part
    $d\alpha$ carries the topologically trivial class $[\,d\alpha\,]=0$ (status
    $\Hfont$). We caution that this is \emph{global} cohomological triviality: at
    the level of a potential $A$ with $\omega=dA$, the genuine local gauge
    redundancy is $A\mapsto A+d\lambda$, which leaves $\omega$ pointwise
    invariant, whereas a nonzero exact field strength $d\alpha$ need not vanish
    pointwise. Thus ``pure gauge'' here means null in cohomology, not $\omega=0$.
  \item[(iii)] For a closed oriented $n$-cycle $\Gamma$, the flux
    $\Phi(\Gamma):=\int_\Gamma\omega$ depends only on $[\omega]$ and the homology
    class $[\Gamma]$, by Stokes' theorem; it is the conserved charge measured over
    $\Gamma$.
\end{enumerate}
\end{proposition}

\begin{proof}
(i)--(ii) are immediate from \Cref{def:derham}: $d\omega=0$ means $\omega$ is a
de~Rham cocycle, and $d(\omega+d\alpha)=d\omega=0$ with
$[\omega+d\alpha]=[\omega]$ since $d\alpha$ is exact. (iii): if $\omega'=\omega+d\alpha$
and $\Gamma'=\Gamma+\partial S$, then by Stokes
$\int_{\Gamma'}\omega'=\int_\Gamma\omega+\int_\Gamma d\alpha+\int_{\partial S}\omega+\int_{\partial S}d\alpha$,
and each correction integral vanishes because $\Gamma$ is closed
($\int_\Gamma d\alpha=\int_{\partial\Gamma}\alpha=0$), $\omega$ is closed
($\int_{\partial S}\omega=\int_S d\omega=0$), and likewise for the last term.
Thus $\Phi$ descends to $H^n_{\mathrm{dR}}\times H_n$.
\end{proof}

By the de~Rham theorem, $H^n_{\mathrm{dR}}(X)\cong H^n(X;\Rm)$ with the pairing of
\Cref{def:pairing} matching integration; so the de~Rham and singular pictures are
the same status-$\Sfont$ observable module in two realization channels.

\subsection{Exact sequences: conservation and obstruction bookkeeping}

\begin{proposition}[Long exact sequence of a pair]
\label{prop:les}
For a pair $(X,B)$ there is a natural long exact sequence
\[
\begin{tikzcd}[column sep=1.6em]
\cdots \arrow[r] & H_n(B) \arrow[r, "i_*"] & H_n(X) \arrow[r, "j_*"] &
H_n(X,B) \arrow[r, "\partial"] & H_{n-1}(B) \arrow[r] & \cdots
\end{tikzcd}
\]
in which $\im i_* = \ker j_*$, $\im j_*=\ker\partial$, and
$\im\partial=\ker i_*$. Physically, $\partial$ is the boundary/defect
charge-transfer map: a bulk class with nontrivial $\partial$ deposits conserved
charge on the boundary $B$.
\end{proposition}

\begin{proof}
This is the homology long exact sequence associated to the short exact sequence of
chain complexes $0\to C_\bullet(B)\to C_\bullet(X)\to C_\bullet(X,B)\to 0$, where
$C_\bullet(X,B):=C_\bullet(X)/C_\bullet(B)$. Exactness is the standard homological
snake lemma / zig-zag construction: the connecting map $\partial$ sends the class
of a relative cycle $c$ (so $\partial c\in C_{n-1}(B)$) to $[\partial c]\in
H_{n-1}(B)$, and the three exactness statements are verified by diagram chase.
Naturality follows because a map of pairs induces a map of the short exact
sequences.
\end{proof}

\begin{proposition}[Mayer--Vietoris: reconstructing global physics from patches]
\label{prop:mv}
If $X=U\cup V$ with $U,V$ open, there is a long exact sequence
\[
\begin{tikzcd}[column sep=1.35em]
\cdots \arrow[r] & H_n(U\cap V) \arrow[r] & H_n(U)\oplus H_n(V) \arrow[r] &
H_n(X) \arrow[r, "\partial"] & H_{n-1}(U\cap V) \arrow[r] & \cdots
\end{tikzcd}
\]
This is the topological form of the Locality/descent axiom
(\Cref{ax:locality}): global conserved information on $X$ is reconstructed from
local information on $U$, $V$ together with the gluing data on the overlap
$U\cap V$; the connecting map $\partial$ is the obstruction to naive gluing.
\end{proposition}

\begin{proof}
Apply the homology long exact sequence to the short exact sequence
$0\to C_\bullet(U\cap V)\xrightarrow{(i,-j)}C_\bullet(U)\oplus C_\bullet(V)
\xrightarrow{k+l}C_\bullet(U+V)\to 0$, and use that the inclusion
$C_\bullet(U+V)\hookrightarrow C_\bullet(X)$ of sums of chains supported in $U$ or
$V$ induces isomorphisms on homology (small-simplices / barycentric-subdivision
theorem). Exactness is again the zig-zag lemma.
\end{proof}

\begin{remark}[Descent reading]
\label{rem:mv-descent}
\Cref{prop:mv} is exactly a descent statement: the sequence exhibits $H_\bullet(X)$
as glued from $H_\bullet(U),H_\bullet(V)$ along $H_\bullet(U\cap V)$, up to the
correction measured by $\partial$. When $\partial=0$ the local data glue on the
nose (a \emph{sheaf}-like situation); when $\partial\ne 0$ they glue only up to
the obstruction (a \emph{stack}-like situation). This is the first appearance in
the ladder of the sheaf/stack dichotomy that Part~VI makes precise.
\end{remark}

\subsection{Poincar\'e duality and cup product}

\begin{theorem}[Poincar\'e duality]
\label{thm:pd}
Let $X$ be a closed oriented $n$-manifold with fundamental class $[X]\in H_n(X)$.
Then cap product with $[X]$ is an isomorphism
\[
  -\frown[X]\colon H^k(X)\xrightarrow{\ \cong\ }H_{n-k}(X)\qquad(0\le k\le n).
\]
Physically: a codimension-$k$ defect (a class in $H_{n-k}$) is equivalent data to
a degree-$k$ field-strength observable (a class in $H^k$); defects and dual field
strengths carry the same conserved information.
\end{theorem}

\begin{proof}
This is the classical Poincar\'e duality theorem; we recall the mechanism. Cap
product $\frown\colon H^k(X)\otimes H_n(X)\to H_{n-k}(X)$ is defined on the chain
level by $\varphi\frown\sigma=\varphi(\sigma\!\mid_{\text{front }k})\,\sigma\!\mid_{\text{back }n-k}$
and descends to (co)homology. For a closed oriented manifold, capping with the
fundamental class is shown to be an isomorphism by the standard local-to-global
argument, which is naturally phrased using \emph{compactly supported} cohomology
$H^k_c$: unlike ordinary cohomology (which has $H^n(\Rm^n)=0$), one has
$H^k_c(\Rm^n)\cong\Zm$ concentrated in degree $k=n$, and the duality
$H^k_c(\Rm^n)\cong H_{n-k}(\Rm^n)$ holds for $\Rm^n$ and half-spaces (where both
sides are computable). One then propagates along a good cover by an induction
using the Mayer--Vietoris sequences (\Cref{prop:mv}) on both sides and the five
lemma; the naturality of $\frown$ makes the induction squares commute. For the
closed manifold $X$ one has $H^k_c(X)=H^k(X)$, yielding the stated form.
\end{proof}

\begin{definition}[Cup product]
\label{def:cup}
The \emph{cup product} $\smile\colon H^k(X;R)\otimes H^\ell(X;R)\to H^{k+\ell}(X;R)$
is induced by $(\varphi\smile\psi)(\sigma)=\varphi(\sigma\!\mid_{\text{front}})\cdot
\psi(\sigma\!\mid_{\text{back}})$, making $H^\bullet(X;R)$ a graded-commutative
ring. Physically the cup product is the \emph{interaction} of two charge/flux
observables, and its graded commutativity $\alpha\smile\beta=(-1)^{|\alpha||\beta|}\beta\smile\alpha$
$(\Sfont/\Hfont)$ encodes the statistics of the corresponding excitations.
\end{definition}

\begin{example}[Conserved charge on the torus]
\label{ex:torus}
For $X=T^2=S^1\times S^1$ one has $H^1(T^2;\Zm)\cong\Zm^2$ generated by classes
$a,b$ dual to the two circles, with $a\smile a=b\smile b=0$ and $a\smile b=-b\smile a$
generating $H^2(T^2;\Zm)\cong\Zm$. Two \emph{distinct} conserved observables live
here and must not be confused.
\begin{enumerate}
  \item[(i)] A \emph{flat} gauge field ($F=0$) is classified up to gauge by its
    two holonomies, i.e.\ by a class in $H^1(T^2;\Umone)\cong\Umone\times\Umone$
    (Wilson loops around the two cycles). Such a connection has vanishing
    curvature and hence \emph{trivial} first Chern number; the moduli space of
    flat connections detects only the $H^1$-data.
  \item[(ii)] A generally \emph{non-flat} complex line bundle $L\to T^2$ carries a
    quantized magnetic flux, the first Chern number
    $\int_{T^2}\tfrac{F}{2\pi}=\langle c_1(L),[T^2]\rangle\in\Zm$ with
    $c_1(L)\in H^2(T^2;\Zm)\cong\Zm$; a nonzero value \emph{forces} $F\neq 0$, so
    this observable is disjoint from case (i).
\end{enumerate}
It is precisely the nondegenerate cup pairing
$\langle a\smile b,[T^2]\rangle=1$ that makes $H^2(T^2;\Zm)\cong\Zm$
one-dimensional and thereby able to host the integer flux of case (ii). Both are
toy models underlying the quantum-Hall and surface-code phenomenology revisited in
Part~VI.
\end{example}

% =====================================================================
\section{Characteristic classes and the index theorem}
\label{sec:char}

Characteristic classes are the module of \emph{quantized} topological information:
they assign to a gauge bundle a cohomology class of its base, natural under
pullback, and their integrals are integer- or rational-valued conserved charges.

\subsection{Classifying spaces and characteristic classes}

\begin{definition}[Principal bundle and classifying space]
\label{def:bundle}
A \emph{principal $G$-bundle} over $X$ is a fiber bundle $\pi\colon P\to X$ with a
free, fiber-transitive right $G$-action. Isomorphism classes of principal
$G$-bundles over a paracompact $X$ are in natural bijection with homotopy classes
of maps $[X,BG]$, where $BG$ is the \emph{classifying space} (the base of the
universal bundle $EG\to BG$); the bijection sends $f\colon X\to BG$ to the
pullback $f^*EG$.
\end{definition}

\begin{definition}[Characteristic class]
\label{def:charclass}
A \emph{characteristic class} for $G$-bundles with coefficients in $R$ is a natural
transformation $c\colon [-,BG]\Rightarrow H^*(-;R)$; equivalently, by
\Cref{def:bundle} and the Yoneda lemma, an element $c\in H^*(BG;R)$, with the
class of a bundle $f^*EG$ given by $c(f^*EG)=f^*c$.
\end{definition}

\begin{proposition}[Naturality is conservation]
\label{prop:char-nat}
For any continuous $g\colon Y\to X$ and any principal $G$-bundle $P\to X$,
$c(g^*P)=g^*c(P)$. In particular characteristic numbers $\int_X c(P)$ of closed
manifolds are homotopy/cobordism invariants: they are conserved under continuous
deformation of the bundle and of the base.
\end{proposition}

\begin{proof}
Naturality is \Cref{def:charclass} directly: if $P=f^*EG$ then
$g^*P=(f\circ g)^*EG$ has class $(f\circ g)^*c=g^*f^*c=g^*c(P)$. The characteristic
number statement follows because $\int_X\colon H^n(X)\to R$ (for $X$ closed
oriented $n$-manifold) is a homotopy invariant of the pair $(X,\text{class})$ by
\Cref{thm:homotopy-inv} and \Cref{prop:coho-functor}.
\end{proof}

\begin{example}[Chern classes and quantized flux]
\label{ex:chern}
For $G=\mathrm{U}(k)$ one has the polynomial ring
\[
  H^*(B\mathrm{U}(k);\Zm)=\Zm[c_1,\dots,c_k],\qquad c_i\in H^{2i},
\]
whose generators $c_i$ are the \emph{Chern classes}. For a complex line bundle
$L\to X$, $c_1(L)\in H^2(X;\Zm)$ is the quantized first Chern number; over a closed
surface $\Sigma$, $\int_\Sigma c_1(L)=\frac{1}{2\pi}\int_\Sigma F\in\Zm$ is the
Dirac-quantized magnetic monopole charge / total flux. The \emph{Chern character}
\[
  \ch(E)=\rk(E)+c_1+\tfrac12(c_1^2-2c_2)+\cdots\in H^{\mathrm{even}}(X;\Qm)
\]
is a ring homomorphism $\ch\colon K^0(X)\to H^{\mathrm{even}}(X;\Qm)$ turning
$K$-theory charges into cohomological ones. All entries here are status-$\Sfont$.
\end{example}

\begin{example}[Stiefel--Whitney classes and spin]
\label{ex:sw}
For $G=\mathrm{O}(k)$, the \emph{Stiefel--Whitney classes}
$w_i\in H^i(X;\Zm/2)$ obstruct structure: $w_1(TX)=0$ iff $X$ is orientable, and
(given $w_1=0$) $w_2(TX)=0$ iff $X$ admits a \emph{spin structure}, the condition
under which fermions can be globally defined. Thus $w_2$ is a conserved $\Zm/2$
obstruction whose vanishing is a consistency condition for fermionic matter
$(\Sfont)$.
\end{example}

\subsection{The index theorem as a realization-pipeline identity}

\begin{theorem}[Atiyah--Singer as $\text{observable}=\Real(\text{abstract structure})$]
\label{thm:index}
Let $X$ be a closed even-dimensional spin manifold with spinor bundle
$S=S^+\oplus S^-$ (the chiral splitting), and $E\to X$ a complex vector bundle
with connection; let $D_E^+\colon\Gamma(S^+\otimes E)\to\Gamma(S^-\otimes E)$ be
the associated chiral twisted Dirac operator. Then the analytic index equals a
purely topological expression:
\begin{equation}
\label{eq:index}
  \ind(D_E^+) \;=\; \dim\ker D_E^+ - \dim\ker D_E^-
  \;=\; \int_X \Ahat(TX)\wedge\ch(E)\;\in\;\Zm ,
\end{equation}
where $D_E^-$ is the formal adjoint (so $\ker D_E^-=\operatorname{coker}D_E^+$).
Reading the left side as an \emph{observable} $\Obs$ (a net chiral zero-mode
count / anomaly coefficient) and the right side as
$\Real_{\mathrm{top}}\bigl(\Phi(X,E)\bigr)$ with $\Phi(X,E)$ the topological data
$\Ahat(TX),\ch(E)$, \eqref{eq:index} is an instance of the Part~I master formula
$\Obs_{\mathrm{top}}(M)=\Obs(\Real_{\mathrm{top}}(\Phi(M)))$ with realization
channel $\alpha=\mathrm{top}$. Status $\Sfont$.
\end{theorem}

\begin{proof}
Equation \eqref{eq:index} is the Atiyah--Singer index theorem
\cite{AtiyahSinger}, which we do not reprove; we verify only the pipeline reading.
Set $\Phi(X,E):=(\Ahat(TX),\ch(E))\in H^\bullet(X;\Qm)$, a functor of the
topological data (natural under bundle pullback by \Cref{prop:char-nat}). Let
$\Real_{\mathrm{top}}$ be the integration channel
$\Real_{\mathrm{top}}(\eta)=\int_X\eta$ applied to the top-degree component of the
product $\Ahat(TX)\wedge\ch(E)$, valued in $\Qm$; and let $\Obs$ be the identity on
$\Zm\subseteq\Qm$. Then $\Obs(\Real_{\mathrm{top}}(\Phi(X,E)))=\int_X\Ahat(TX)\wedge\ch(E)$,
which by \eqref{eq:index} equals $\ind(D_E^+)$, the analytic observable. Each of
$\Phi,\Real_{\mathrm{top}},\Obs$ is functorial (naturality, linearity of the
integral, inclusion), so $\Obs_{\mathrm{top}}=\Obs\circ\Real_{\mathrm{top}}\circ\Phi$
is a realization pipeline in the sense of \Cref{def:pipeline}, and \eqref{eq:index}
is the assertion that it computes the index. Integrality of the right side is a
nontrivial consequence, exactly the content of the theorem.
\end{proof}

\begin{remark}[Why the index theorem is the ideal status-$\Sfont$ anchor]
\label{rem:index-anchor}
\Cref{thm:index} imports no new mathematics: its physical reading
(``topology $\to$ analysis: zero modes and anomalies from topology'') \emph{is} the
theorem. This makes it the least contestable entry in the entire library and the
model we hold other, more heuristic entries against.
\end{remark}

% =====================================================================
\section{Topological quantum field theory: functorial realization}
\label{sec:tqft}

We now reach the conceptual center of Part~IV. A TQFT is a realization pipeline
that is \emph{itself a functor} over a category of cobordisms; it is the first
status-$\Sfont$ instance of the Part~I pipeline in which $\Real_\alpha$ is a
genuine functor, not a heuristic translation rule. This is precisely why Part~IV
is where category theory enters the ladder as ambient grammar.

\subsection{The bordism category}

\begin{definition}[Bordism category]
\label{def:bord}
Fix $n\ge 1$. The \emph{bordism category} $\Bord_n$ has:
\begin{itemize}
  \item \textbf{objects}: closed oriented $(n-1)$-manifolds $\Sigma$;
  \item \textbf{morphisms} $\Sigma_0\to\Sigma_1$: oriented cobordisms, i.e.\
    (diffeomorphism classes rel boundary of) compact oriented $n$-manifolds $W$
    with $\partial W\cong \overline{\Sigma_0}\sqcup\Sigma_1$;
  \item \textbf{composition}: gluing of cobordisms along the common boundary;
  \item \textbf{identities}: cylinders $\Sigma\times[0,1]$.
\end{itemize}
Disjoint union $\sqcup$ makes $\Bord_n$ a symmetric monoidal category with unit the
empty $(n-1)$-manifold $\varnothing$.
\end{definition}

\begin{definition}[TQFT, Atiyah's axioms]
\label{def:tqft}
An \emph{$n$-dimensional TQFT} valued in $\Vect_k$ (finite-dimensional
$k$-vector spaces) is a symmetric monoidal functor $Z\colon\Bord_n\to\Vect_k$.
Explicitly:
\begin{enumerate}
  \item[(Z1)] $Z$ assigns to each closed $(n-1)$-manifold $\Sigma$ a
    finite-dimensional vector space $Z(\Sigma)$, with $Z(\varnothing)=k$;
  \item[(Z2)] $Z(\Sigma_0\sqcup\Sigma_1)\cong Z(\Sigma_0)\otimes Z(\Sigma_1)$
    (monoidality);
  \item[(Z3)] $Z$ assigns to a cobordism $W\colon\Sigma_0\to\Sigma_1$ a linear map
    $Z(W)\colon Z(\Sigma_0)\to Z(\Sigma_1)$, with
    $Z(W'\circ W)=Z(W')\circ Z(W)$ (gluing) and $Z(\Sigma\times[0,1])=\id_{Z(\Sigma)}$;
  \item[(Z4)] $Z(\overline\Sigma)\cong Z(\Sigma)^*$ (orientation reversal is dual).
\end{enumerate}
A closed $n$-manifold $M$ is a cobordism $\varnothing\to\varnothing$, so
$Z(M)\in\Hom(k,k)=k$ is a numerical invariant (a partition function).
\end{definition}

\subsection{TQFT is a representation entry}

\begin{theorem}[TQFT realizes the Part~I pipeline; monoidality is the Decomposition axiom]
\label{thm:tqft-entry}
An $n$-dimensional TQFT $Z\colon\Bord_n\to\Vect_k$ is precisely a status-$\Sfont$
representation entry $E=(\Bord_n,\Vect_k,Z,\Sfont)$ in the sense of
\Cref{def:entry}, realizing the pipeline of \Cref{def:pipeline} with
\[
\begin{gathered}
  \Phi=\text{(inclusion of a closed $(n-1)$-manifold as an object)},\\
  \Real_\alpha=Z,\qquad
  \Obs=\dim\ (\text{or trace/partition function}).
\end{gathered}
\]
Moreover the symmetric monoidal structure of $Z$ is exactly the Decomposition
axiom (\Cref{ax:decomp}) applied to disjoint unions and gluing:
\begin{equation}
\label{eq:tqft-decomp}
  Z(M_1\sqcup M_2)=Z(M_1)\otimes Z(M_2),\qquad
  Z(W_2\circ W_1)=Z(W_2)\circ Z(W_1).
\end{equation}
\end{theorem}

\begin{proof}
The data of \Cref{def:entry} are: a mathematical structure $M=\Bord_n\in\Math_{\Man}$,
a physical target $P=\Vect_k\in\Phys$ (state spaces and linear processes), a
translation datum $\tau=Z$, and a status $\sigma$. Since $Z$ is an
\emph{actual functor} (Atiyah's axioms), $\tau$ is functorial and the entry is
status-$\Sfont$ by \Cref{rem:framework-status}. To exhibit the pipeline, note
$\Phi$ includes an $(n-1)$-manifold $\Sigma$ as an object of $\Bord_n$;
$\Real_\alpha=Z$ sends it to $Z(\Sigma)\in\Vect_k$; and $\Obs=\dim$ (or, on closed
$M$, the scalar $Z(M)$) extracts a number. Each is functorial and they compose,
so $\Obs_\alpha=\Obs\circ Z\circ\Phi$ is a realization pipeline
(\Cref{def:pipeline}).

Finally, \eqref{eq:tqft-decomp} is the statement that $Z$ is symmetric monoidal
(Z2) and a functor (Z3). Reading disjoint union as ``parallel composition of
systems'' and gluing as ``sequential composition of processes,'' \eqref{eq:tqft-decomp}
says the observable of a composite decomposes into the observables of its pieces
via $\otimes$ and $\circ$; this is precisely the Decomposition axiom instantiated
on cobordisms. Uniqueness of the entry follows because a symmetric monoidal
functor is determined by its action on objects and generating morphisms, which is
the data specified.
\end{proof}

\begin{example}[A 2d TQFT is a commutative Frobenius algebra]
\label{ex:2d-tqft}
For $n=2$, the classification theorem states that $2$-dimensional TQFTs
$Z\colon\Bord_2\to\Vect_k$ correspond bijectively to commutative Frobenius
algebras $A=Z(S^1)$: the pair of pants (a cobordism $S^1\sqcup S^1\to S^1$) gives
the multiplication $A\otimes A\to A$, the reversed pair of pants gives the
comultiplication, the disk ($\varnothing\to S^1$) gives the unit
$\eta\colon k\to A$, and the reversed disk ($S^1\to\varnothing$) gives the counit
$\varepsilon\colon A\to k$ (the Frobenius trace form). Gluing the two disks yields
the sphere $S^2$ ($\varnothing\to\varnothing$), which therefore evaluates to the
scalar $\varepsilon(\eta(1))=\varepsilon(1_A)$; the closed torus $T^2$, obtained
instead as the trace of $\id_A$ (composition of coevaluation and evaluation),
evaluates to $\dim_k A$. Associativity, commutativity, and the
Frobenius relation are
each a diffeomorphism between glued surfaces, hence forced by functoriality
\eqref{eq:tqft-decomp}. This is the cleanest illustration that ``the algebra of
observables $=$ the decomposition law of cobordisms.'' Status $\Sfont$.
\end{example}

\subsection{The cobordism hypothesis}

\begin{theorem}[Cobordism hypothesis; Baez--Dolan, Lurie]
\label{thm:cobordism-hyp}
Let $\mathcal{C}$ be a symmetric monoidal $(\infty,n)$-category with duals. The
evaluation-at-a-point map
\[
  \operatorname{Fun}^{\otimes}\bigl(\Bord_n^{\mathrm{fr}},\mathcal{C}\bigr)
  \xrightarrow{\ Z\mapsto Z(\mathrm{pt})\ }
  \mathcal{C}\sfd
\]
is an equivalence between the space of fully extended framed $n$-dimensional TQFTs
valued in $\mathcal{C}$ and the maximal $\infty$-groupoid $\mathcal{C}\sfd$ of
\emph{fully dualizable} objects of $\mathcal{C}$. Thus a fully extended TQFT is
determined, up to a contractible space of choices, by a single fully dualizable
object $Z(\mathrm{pt})$. Status $\Sfont$ (theorem of Lurie; proof sketched at
length in \cite{Lurie2009}).
\end{theorem}

\begin{proof}[Proof idea (cited, not reproduced)]
The full proof \cite{Lurie2009} is an induction on codimension. The key input is
that $\Bord_n^{\mathrm{fr}}$ is the \emph{free} symmetric monoidal
$(\infty,n)$-category with duals on a single object $\mathrm{pt}$: every
cobordism is generated, under composition and duality, by the point together with
its (higher) evaluation and coevaluation morphisms. A symmetric monoidal functor
out of a free such category is determined by the image of the generator, subject
to that image being fully dualizable (so that all the duality data have targets).
Hence $Z\mapsto Z(\mathrm{pt})$ is an equivalence onto $\mathcal{C}\sfd$.
Framings can be traded for tangential $G$-structures by taking homotopy fixed
points of an $O(n)$-action, giving the structured variants.
\end{proof}

\begin{corollary}[Rigidity of topological realization channels]
\label{cor:rigidity}
In the extended setting, the topological realization channel is \emph{rigid}: once
the point-realization $Z(\mathrm{pt})$ is fixed as a fully dualizable object, all
higher-codimension observables (line, surface, and defect data) are uniquely
determined. Physically, fixing how the elementary local excitation realizes fixes
the entire extended field theory. This is the strongest compositionality statement
in the ladder and directly motivates Part~V's abstraction of $\mathcal{C}$ into the
general theory of (higher) categories.
\end{corollary}

\begin{proof}
Immediate from \Cref{thm:cobordism-hyp}: the evaluation map is an equivalence, so
its inverse reconstructs $Z$ (all its values in every codimension) from
$Z(\mathrm{pt})$ functorially and essentially uniquely.
\end{proof}

The functorial picture is summarized by the commuting triangle of monoidal
functors, in which $Z$ factors the Part~I pipeline through the topological
channel:
\[
\begin{tikzcd}[column sep=3.2em, row sep=large]
\Bord_n \arrow[r, "Z"] \arrow[dr, "\Obs_\alpha"'] & \Vect_k \arrow[d, "\dim\,/\,\mathrm{tr}"]\\
& k
\end{tikzcd}
\qquad\qquad
\begin{tikzcd}[column sep=2.4em]
\mathrm{pt} \arrow[r, mapsto, "Z"] & Z(\mathrm{pt})\in\mathcal{C}\sfd
\end{tikzcd}
\]

% =====================================================================
\section{Dijkgraaf--Witten theory: locality from group cohomology}
\label{sec:dw}

Dijkgraaf--Witten (DW) theory is a concrete, computable TQFT built from a finite
group and a group cocycle; it is the ideal illustration of the Locality/descent
axiom, and of the coboundary $=$ gauge reading of the dictionary.

\begin{definition}[Group cohomology and DW data]
\label{def:dw-data}
For a finite group $G$ and coefficients $\Rm/\Zm$ (trivial action), the
\emph{group cohomology} $H^\bullet(G;\Rm/\Zm):=H^\bullet(BG;\Rm/\Zm)$ is computed
by the bar complex: $n$-cochains are functions $\varphi\colon G^n\to\Rm/\Zm$, with
\emph{additive} coboundary
\begin{multline*}
  (\delta\varphi)(g_1,\dots,g_{n+1})=\varphi(g_2,\dots,g_{n+1})\\
  +\sum_{i=1}^{n}(-1)^i\,\varphi(g_1,\dots,g_ig_{i+1},\dots,g_{n+1})
  +(-1)^{n+1}\varphi(g_1,\dots,g_n) .
\end{multline*}
Here the leading term $\varphi(g_2,\dots,g_{n+1})$ carries no prefactor precisely
because the $G$-action on the coefficients $\Rm/\Zm$ is \emph{trivial}; for a
nontrivial action it would instead read $g_1\!\cdot\!\varphi(g_2,\dots,g_{n+1})$.
A \emph{Dijkgraaf--Witten datum} in dimension $3$ is a class
$\omega\in H^3(G;\Rm/\Zm)$; the associated phase $e^{2\pi i\,\omega}\in\Umone$
recovers the equivalent multiplicative $\Umone$-valued convention.
\end{definition}

\begin{theorem}[DW theory instantiates the Locality/descent axiom]
\label{thm:dw}
Let $G$ be finite and $\omega\in H^3(BG;\Rm/\Zm)$ a $3$-cocycle. The state sum
\begin{equation}
\label{eq:dw}
  Z_\omega(X)=\frac{1}{|G|}\sum_{\phi\in\Hom(\pi_1(X),G)}
  \exp\!\bigl(2\pi i\,\langle\phi^*\omega,[X]\rangle\bigr)
\end{equation}
over a closed oriented $3$-manifold $X$ defines a TQFT
$Z_\omega\colon\Bord_3\to\Vect_\Cm$; here the pairing
$\langle\phi^*\omega,[X]\rangle$ lies in $\Rm/\Zm$, so the exponential
$e^{2\pi i(-)}\in\Umone$ is well typed. The state sum is independent of the
triangulation (hence well defined and local in the sense of
\Cref{ax:locality}) precisely because $\omega$ is a cocycle; and a coboundary shift
$\omega\mapsto\omega+\delta\eta$ changes $Z_\omega$ by nothing, matching the
``coboundary $=$ pure gauge'' dictionary entry.
\end{theorem}

\begin{proof}
Present $X$ by a triangulation $T$ with an ordering of vertices; a discrete
$G$-connection is a labelling of edges by $G$ compatible on faces, i.e.\ a map
$\phi\colon\pi_1(X)\to G$ up to gauge. Each tetrahedron contributes a phase
$\exp\!\bigl(\pm 2\pi i\,\omega(g,h,k)\bigr)$ according to its orientation, and the
state sum \eqref{eq:dw} is the product of these phases over tetrahedra, averaged
over connections. Two triangulations of $X$ are related by a finite sequence of
\emph{Pachner (bistellar) moves}. The invariance of the local weight under each
Pachner move is \emph{equivalent} to the cocycle condition $\delta\omega=0$: the
$2$--$3$ Pachner move is exactly the (additive) pentagon/cocycle identity
\[
  \omega(h,k,l)-\omega(gh,k,l)+\omega(g,hk,l)-\omega(g,h,kl)+\omega(g,h,k)=0
  \quad\text{in }\Rm/\Zm .
\]
Hence $Z_\omega(X)$ is independent of $T$, establishing locality and gluing. If
$\omega'=\omega+\delta\eta$ for a $2$-cochain $\eta$, the extra phases are
associated to faces and their exponents telescope to $0$ (mod $\Zm$) over the
closed manifold $X$ (each internal face is shared by two tetrahedra with opposite
orientation), so $Z_{\omega'}(X)=Z_\omega(X)$. Functoriality on cobordisms follows
because the
state sum on a manifold with boundary defines a vector in the state space of the
boundary and gluing corresponds to contracting these vectors.
\end{proof}

\begin{remark}[Descent and the Equivalence axiom]
\label{rem:dw-equiv}
\Cref{thm:dw} realizes the dashed arrow of the Part~I master diagram in a discrete
setting: the equivalence $\omega\sim\omega+\delta\eta$ (change of cocycle
representative) is a gauge/equivalence relation under which the physical content
$Z_\omega$ is invariant. Thus DW theory simultaneously discharges the
Locality/descent axiom (triangulation independence) and the Equivalence axiom
(coboundary invariance). Status $\Sfont$.
\end{remark}

\begin{example}[$G=\Zm/N$, untwisted]
\label{ex:zn-dw}
For $\omega=0$, $Z_0(X)=|\Hom(\pi_1(X),\Zm/N)|/N=|H^1(X;\Zm/N)|/N$. On the
$3$-torus $T^3$ this counts flat $\Zm/N$ connections, $|H^1(T^3;\Zm/N)|=N^3$, so
$Z_0(T^3)=N^2$; the ground-state degeneracy $N^2$ is the physical signature of
$\Zm/N$ topological order on the spatial torus, a status-$\Sfont$ topological
invariant recovered directly from $H^1$.
\end{example}

% =====================================================================
\section{Bordism, anomalies, and symmetry-protected topological phases}
\label{sec:bordism}

The final layer of Part~IV classifies \emph{invertible} topological field
theories --- the mathematical home of anomalies and symmetry-protected
topological (SPT) phases --- by bordism-theoretic invariants.

\begin{definition}[Bordism group]
\label{def:bordism}
For a tangential structure $H$ (e.g.\ orientation, spin, or a map to $BG$), the
\emph{bordism group} $\Omega_n^H$ is the abelian group of closed $n$-manifolds
with $H$-structure modulo the relation $M\sim M'$ iff $M\sqcup\overline{M'}=\partial W$
for some compact $H$-manifold $W$; addition is disjoint union. By the
Pontryagin--Thom construction, $\Omega_n^H\cong\pi_n(MTH)$, the $n$-th stable
homotopy group of the Thom spectrum $MTH$.
\end{definition}

\begin{definition}[Invertible field theory]
\label{def:invertible}
A field theory $Z$ is \emph{invertible} if $Z(\Sigma)$ is invertible in the target
symmetric monoidal category for every $\Sigma$ (e.g.\ a line in $\Vect$) and each
$Z(W)$ is an isomorphism. Invertible TQFTs are exactly the fully dualizable
objects with trivial higher morphisms, hence factor through a spectrum-level map
by \Cref{thm:cobordism-hyp}.
\end{definition}

\begin{theorem}[Freed--Hopkins classification of invertible theories; SPT phases]
\label{thm:spt}
Reflection-positive invertible $n$-dimensional field theories with tangential
structure $H$ are classified by the torsion subgroup of
\[
  \bigl[MTH,\ \Sigma^{n+1}I\Zm\bigr] \;\cong\; \bigl(I\Zm\bigr)^{n+1}(MTH),
\]
the degree-$(n+1)$ Anderson-dual cohomology of the Thom spectrum $MTH$
(recall $E^k(X)=[X,\Sigma^k E]$, so an $(n+1)$-fold suspension lands in degree
$n+1$)
\cite{FreedHopkins}. Consequently, a symmetry-protected topological phase with
finite symmetry group $G$ in $d$ spacetime dimensions is classified, at the level
of the relevant (co)bordism invariant, by a summand of
$\Omega^{H}_{d}(BG)$ or its generalized-cohomology refinement. The pure-mathematics
classification is status $\Sfont$; the identification of a \emph{specific} lattice
model's phase with a given bordism invariant is model-dependent, status
$\Sfont/\Hfont$.
\end{theorem}

\begin{proof}[Proof idea (cited)]
By \Cref{def:invertible} an invertible theory is a map of spectra from the
bordism spectrum $MTH$ to the target; reflection positivity forces the target to
be the Anderson dual $I\Zm$ of the sphere (Freed--Hopkins \cite{FreedHopkins},
building on \cite{Freed2014}). The classification is then the computation of the
homotopy classes $[MTH,\Sigma^{n+1}I\Zm]$, which by definition of Anderson duality
is the stated generalized cohomology group. Coupling to a background $G$-symmetry
replaces $MTH$ by $MTH\wedge BG_+$, giving the $\Omega^H_d(BG)$-level statement.
The physical identification with an SPT phase is via the partition function on
mapping tori and its response to symmetry defects, which is model-dependent, hence
the composite label.
\end{proof}

\begin{remark}[Anomaly as failure of descent]
\label{rem:anomaly}
An anomaly of an $n$-dimensional theory is precisely a nontrivial invertible
$(n+1)$-dimensional theory in which it lives as a boundary --- a failure of the
partition function to be a number rather than a vector, i.e.\ a failure of the
Locality/descent axiom (\Cref{ax:locality}) at the top codimension. Bordism
invariants thus \emph{measure} the obstruction to descent, closing the loop with
\Cref{rem:mv-descent}: anomalies are the top-dimensional analogue of the
Mayer--Vietoris connecting map. Status $\Sfont/\Hfont$.
\end{remark}

% =====================================================================
\section{Results}
\label{sec:results}

We collect the principal results of Part~IV and their status labels.

\begin{itemize}
  \item \textbf{Homotopy invariance (\Cref{thm:homotopy-inv})} and the
    \textbf{conserved global information functor} $H_\bullet\colon\hTop\to\GrAb$
    (\Cref{def:cgi}): topological observables are conserved under continuous
    deformation. Status $\Sfont$; physical reading $\Sfont/\Hfont$.
  \item \textbf{Conservation principle (\Cref{thm:cons-functor})}: over a
    contractible parameter base, (co)homology classes and their pairings are
    constant. Status $\Sfont$.
  \item \textbf{de~Rham gauge picture (\Cref{prop:derham-gauge})} and
    \textbf{Poincar\'e duality (\Cref{thm:pd})}: closed forms are conserved
    fields, exact forms pure gauge, defects dual to field strengths. Status
    $\Sfont$; readings $\Sfont/\Hfont$.
  \item \textbf{Index theorem as pipeline identity (\Cref{thm:index})}: the
    strongest status-$\Sfont$ instance of $\text{observable}=\Real(\text{abstract
    structure})$.
  \item \textbf{TQFT as representation entry (\Cref{thm:tqft-entry})}: the first
    functorial, status-$\Sfont$ realization of the Part~I pipeline; monoidality is
    the Decomposition axiom.
  \item \textbf{Cobordism hypothesis (\Cref{thm:cobordism-hyp}, \Cref{cor:rigidity})}:
    topological realization channels are rigid; the bridge to Part~V. Status
    $\Sfont$.
  \item \textbf{Dijkgraaf--Witten (\Cref{thm:dw})}: Locality/descent and
    Equivalence axioms from a group cocycle. Status $\Sfont$.
  \item \textbf{Bordism classification (\Cref{thm:spt})}: anomalies and SPT phases
    as bordism invariants; anomaly as failure of descent (\Cref{rem:anomaly}).
    Status $\Sfont$ / $\Sfont\!/\Hfont$.
\end{itemize}

\begin{proposition}[Status statistics of Part~IV]
\label{prop:status-stats}
Of the eight headline results above, seven are status-$\Sfont$ as pure
mathematics (all but the model-specific half of \Cref{thm:spt}); the heuristic
content is confined to the \emph{physical readings} of dictionary entries, and no
result is status-$\Pfont$. Part~IV is therefore the most rigorously anchored
module of the ladder to date.
\end{proposition}

\begin{proof}
By inspection of the individual status labels assigned at each theorem: TQFT
functoriality (\Cref{thm:tqft-entry}), the cobordism hypothesis
(\Cref{thm:cobordism-hyp}), the index theorem (\Cref{thm:index}), Dijkgraaf--Witten
theory (\Cref{thm:dw}), homotopy invariance (\Cref{thm:homotopy-inv}), the
conservation principle (\Cref{thm:cons-functor}), and Poincar\'e duality
(\Cref{thm:pd}) are each established theorems of algebraic topology or mathematical
physics; only the model-identification half of \Cref{thm:spt} is $\Sfont/\Hfont$,
and the ontological framing (inherited, not asserted here) is the only place
$\Pfont$ occurs in the parent program. No headline result of Part~IV is
independently $\Pfont$.
\end{proof}

% =====================================================================
\section{Companion formalization: Haskell and Lean}
\label{sec:code}

Following the parent program's formal-verification bridge (``math is code, code is
math''; Curry--Howard--Lambek), Part~IV ships companion code. The Haskell package
compiles under GHC and encodes the structures of \Cref{sec:homology,sec:char,sec:tqft};
the Lean file is a best-effort sketch of the same signatures in a proof-assistant
idiom.

\subsection{Haskell: chain complexes, homology, and TQFT}

The module \texttt{ChainComplex} realizes \Cref{def:chain,def:homology} over
$\Zm$-free chain groups with an explicit boundary matrix, computes Betti numbers
via ranks (Smith-normal-form-free rational homology), and verifies
$\partial^2=0$. The module \texttt{TQFT} encodes \Cref{def:tqft} as a
symmetric-monoidal-functor typeclass and instantiates the $2$d Frobenius-algebra
example (\Cref{ex:2d-tqft}) and the untwisted Dijkgraaf--Witten count
(\Cref{ex:zn-dw}). The following excerpt shows the core signatures:
{\small
\begin{verbatim}
data ChainComplex = CC { dims :: [Int], boundary :: [Matrix] }
d2isZero     :: ChainComplex -> Bool
bettiNumbers :: ChainComplex -> [Int]
class SymMonFunctor z where
  fUnit    :: z
  fTensor  :: z -> z -> z
  fCompose :: z -> z -> z
dwUntwisted  :: Int -> Int -> Double   -- (N, b_1 of manifold) -> Z(T^3)
\end{verbatim}
}
The \texttt{Main} module runs all demonstrations: it verifies $\partial^2=0$ and
prints Betti numbers for the circle, torus, and $2$-sphere (recovering
$b_\bullet(S^1)=(1,1)$, $b_\bullet(T^2)=(1,2,1)$, $b_\bullet(S^2)=(1,0,1)$), checks
the Frobenius/TQFT gluing laws, and prints the Dijkgraaf--Witten degeneracy
$Z_0(T^3)=N^2$ of \Cref{ex:zn-dw}.

\subsection{Lean: chain-complex and homology types}

The Lean sketch \texttt{AlgTop.lean} declares a \texttt{ChainComplex} structure
with fields \texttt{d} (the boundary map in each degree) and a hypothesis
\texttt{d\_comp\_d} asserting $\partial_n\circ\partial_{n+1}=0$ for all $n$,
defines cycles/boundaries/homology as subquotients, and states signatures for the
\texttt{TQFT} structure (a symmetric monoidal functor $\Bord_n\to\Vect$) and for
\Cref{thm:index} and \Cref{thm:dw} as \texttt{theorem} statements with
\texttt{sorry} placeholders. It is best-effort and not build-gated; its purpose is
to record the intended machine-checkable interface, extending Part~I's Lean
skeleton to the topological domain.

% =====================================================================
\section{Discussion}
\label{sec:discussion}

\subsection{How Part~IV builds on Part~III}

Part~III supplies the geometric ``spaces of physical possibility'' (varieties,
moduli, positive geometries). Part~IV attaches to each such space its conserved
global information: (co)homology, characteristic classes, and (via
compactification and degeneration) a bordism class. The boundary stratification of a
positive geometry --- ``boundaries of boundaries'' --- is literally a chain complex
(\Cref{def:chain}), and its homology measures which boundary charges are conserved.
The variation-of-Hodge-structure and Gauss--Manin data of Part~III are the smooth
(de~Rham) realization channel of the same cohomology treated here abstractly, so
Parts~III and~IV are two realization channels ($\alpha=\text{Hodge/de~Rham}$ vs.\
$\alpha=\text{singular/topological}$) of one observable module, matched by the
de~Rham theorem.

\subsection{How Part~IV seeds Part~V}

TQFT (\Cref{thm:tqft-entry}) is the first fully rigorous, status-$\Sfont$ instance
of the Part~I pipeline in which the realization $\Real_\alpha$ is a genuine
functor. The cobordism hypothesis (\Cref{thm:cobordism-hyp}) then shows the entire
extended theory is generated by a single fully dualizable object in a symmetric
monoidal $(\infty,n)$-category $\mathcal{C}$. This forces the ambient grammar of
the next module to be exactly the theory of (higher) categories with duality:
Part~V abstracts $\mathcal{C}$, functoriality, and the duality (evaluation /
coevaluation) data into the general theory of categories, monoidal and
dagger-compact structure, and homotopy type theory. In particular the ``no natural
diagonal'' phenomenon behind no-cloning (Part~V/VI) is the categorical shadow of
the duality data that already organize $\Bord_n$ here.

\subsection{Limitations}

We reproduce and address the parent program's limitations as they bear on Part~IV.
\begin{enumerate}
  \item \emph{Not every analogy is a theory.} The dictionary entries of
    \Cref{sec:dict} carry status labels; the $\Hfont$ entries (e.g.\ ``Postnikov
    tower $=$ hierarchy of defects'') are heuristic and are flagged as such, not
    asserted as established physics.
  \item \emph{No single universal functor $\Math\to\Phys$.} Part~IV constructs
    \emph{domain-local} functors ($H_\bullet$, $Z$), not one global functor; this
    is deliberate and consistent with the modular design.
  \item \emph{Model dependence of SPT identification.} \Cref{thm:spt} is
    status-$\Sfont$ as a classification but only $\Sfont/\Hfont$ when tied to a
    specific lattice model; each such claim needs its own citation.
  \item \emph{Extended/framed subtleties.} \Cref{thm:cobordism-hyp} is stated for
    framed bordism; structured (oriented, spin, $G$-) variants require homotopy
    fixed points of an $O(n)$-action and are more delicate. We flag but do not
    resolve these.
\end{enumerate}

\subsection{Connections to the synthesis backbone}

Three items recur across modules and are load-bearing for the synthesis paper:
(i) the realization pipeline \eqref{eq:pipeline}, instantiated here as $Z$ and as
the index theorem; (ii) homology $H_n$ as conserved cycles, which reappears in
Part~VI as surface-code logical operators $H_1(\Sigma;\Zm/2)$; and (iii) the
sheaf/stack (descent) dichotomy, which appears here as the Mayer--Vietoris
connecting map and the anomaly-as-failed-descent principle (\Cref{rem:mv-descent,rem:anomaly})
and is made rigorous only in Part~VI.

% =====================================================================
\section{Conclusion}
\label{sec:conclusion}

Part~IV formalizes algebraic topology as the grammar of conserved global
information within the Math$\to$Physics representation library. We proved homotopy
invariance and functoriality of homology and packaged them as a conserved global
information functor $H_\bullet\colon\hTop\to\GrAb$; we read closed/exact forms as
conserved fields/pure gauge and defects as dual field strengths; we reframed the
Atiyah--Singer index theorem as an instance of
$\text{observable}=\Real(\text{abstract structure})$; and we exhibited TQFT as the
first functorial, status-$\Sfont$ realization of the Part~I pipeline, with the
cobordism hypothesis as its rigidity/classification capstone and Dijkgraaf--Witten
theory as a computable model of the Locality and Equivalence axioms. The bordism
classification of invertible theories placed anomalies and SPT phases into the same
descent-theoretic frame.

The decisive structural fact is that Part~IV is where \emph{functoriality becomes
achievable in a genuinely topological, status-$\Sfont$ setting}. Every later
appeal to ``the realization functor'' in the ladder is anchored by the concrete
functor $Z\colon\Bord_n\to\Vect$ constructed here. This is precisely the bridge to
Part~V, which abstracts functoriality, monoidal structure, and duality into the
universal grammar of categories and homotopy type theory. In the four-slogan
compression of the parent program --- mathematics is syntax, motives are semantics,
periods are measurements, coalgebras are decomposition laws --- Part~IV supplies the
missing structural clause: \emph{topology is conservation, and functoriality is its
law of composition}.

% =====================================================================
\section*{Acknowledgments}
This work is Part~IV of a seven-part modular program by the YonedaAI Research
Collective. We thank the automated peer-review and formatting-review pipelines
(external reviewers) for iterative feedback incorporated into this version.

% =====================================================================
\appendix
\section{Algebraic topology dictionary (with status labels)}
\label{sec:dict}

The following reproduces the algebraic-topology library with its epistemic status
labels, verbatim in structure from the parent program's Appendix~D. Composite
labels are preserved.

{\small
\begin{longtable}{%
  >{\raggedright\arraybackslash}p{4.2cm}%
  >{\raggedright\arraybackslash}p{4.8cm}%
  >{\raggedright\arraybackslash}p{3.4cm}%
  c}
\toprule
\textbf{Algebraic topology} & \textbf{Physical representation} & \textbf{Semantic meaning} & \textbf{Status}\\
\midrule
\endfirsthead
\multicolumn{4}{l}{\emph{Algebraic topology dictionary (continued)}}\\
\toprule
\textbf{Algebraic topology} & \textbf{Physical representation} & \textbf{Semantic meaning} & \textbf{Status}\\
\midrule
\endhead
\midrule
\multicolumn{4}{r}{\emph{continued on next page}}\\
\endfoot
\bottomrule
\endlastfoot
Topological space $X$ & Configuration, spacetime, or state space & Domain of possibility & $\Sfont$\\
Point & State or event & Localized datum & $\Sfont$\\
Path & Evolution/worldline & Continuous transition & $\Sfont$\\
Loop & Closed evolution & Holonomy probe & $\Sfont$\\
Homotopy & Continuous deformation & Physically equivalent path & $\Sfont$\\
Homotopy class & Topological sector & Deformation-invariant class & $\Sfont$\\
$\pi_1(X)$ & Winding structure & Monodromy, defect winding & $\Sfont$\\
$\pi_n(X)$ & Higher holes & Brane/defect charge sectors & $\Sfont/\Hfont$\\
Homology $H_n(X)$ & Cycles or holes & Conserved extended structures & $\Sfont$\\
Cohomology $H^n(X)$ & Fields/flux observables & Charge detectors on cycles & $\Sfont$\\
Boundary operator $\partial$ & Boundary of region & Where conservation appears & $\Sfont$\\
Cycle $\partial c=0$ & Closed conserved object & Conserved charge carrier & $\Sfont$\\
Boundary $c=\partial b$ & Trivial conserved object & Pure gauge / null sector & $\Hfont$\\
Cocycle & Closed field/charge rule & Conservation condition & $\Sfont$\\
Coboundary & Gauge-exact field & Redundant potential & $\Sfont$\\
Long exact sequence & Boundary--bulk relation & Charge transfer & $\Sfont$\\
Mayer--Vietoris & Gluing local regions & Global from patches & $\Sfont$\\
Cup product & Product of classes & Interaction of observables & $\Sfont/\Hfont$\\
Cap product & Class $\frown$ cycle & Measurement over region & $\Sfont$\\
Poincar\'e duality & Cycles to fields & Defects as dual field strengths & $\Sfont$\\
de~Rham cohomology & Forms mod exact & Gauge-invariant flux & $\Sfont$\\
Closed form & Conserved field & $d\omega=0$ conservation & $\Sfont$\\
Exact form & Gauge-trivial field & $\omega=d\alpha$ redundancy & $\Sfont$\\
Hodge star $\star$ & Dual field operation & EM/geometric duality & $\Sfont$\\
Fiber bundle & Local d.o.f.\ over base & Internal state per point & $\Sfont$\\
Principal bundle & Gauge field topology & Gauge choices glued & $\Sfont$\\
Classifying space $BG$ & Universal classifier & $G$-sector space & $\Sfont$\\
Characteristic class & Bundle invariant & Quantized charge/anomaly & $\Sfont$\\
Chern class & Complex bundle invariant & Monopole, quantized flux & $\Sfont$\\
Chern character & Charge map & Topological response/index & $\Sfont$\\
Stiefel--Whitney class & Real bundle obstruction & Orientation/spin obstruction & $\Sfont$\\
Spin structure & Lift of frame data & Fermion consistency & $\Sfont$\\
Index theorem & Topology to analysis & Zero modes/anomalies & $\Sfont$\\
$K$-theory & Cohomology of bundles & D-brane/phase charges & $\Sfont$\\
Spectra & Stable homotopy objects & Generalized charges & $\Sfont$\\
Generalized cohomology & Beyond ordinary charge & Exotic charges/phases & $\Sfont$\\
Steenrod operations & Cohomology operations & Constraints on phases & $\Sfont/\Hfont$\\
Postnikov tower & Layered homotopy data & Hierarchy of defects & $\Hfont$\\
Loop space $\Omega X$ & Space of closed paths & Recurrence dynamics & $\Sfont/\Hfont$\\
Suspension $\Sigma X$ & Dimension shift & Dimensional lift of defects & $\Hfont$\\
Cobordism & Manifold as transition & Spacetime history & $\Sfont$\\
Bordism group & Manifolds mod boundary & Classification of phases & $\Sfont$\\
TQFT & Functor $\Bord_n\to\Vect$ & Topology-to-algebra & $\Sfont$\\
Extended TQFT & Data in all codimensions & Particles, lines, surfaces & $\Sfont$\\
Fully dualizable object & Point-data of a TQFT & Local seed of global theory & $\Sfont$\\
Group cohomology & Cohomology of symmetry & Discrete action / SPT data & $\Sfont$\\
Dijkgraaf--Witten & Finite-group gauge theory & Action from group cocycle & $\Sfont$\\
\end{longtable}
}

% =====================================================================
\begin{thebibliography}{99}

\bibitem{Atiyah1988}
M.~F. Atiyah,
\emph{Topological quantum field theories},
Publications Math\'ematiques de l'IH\'ES \textbf{68} (1988), 175--186.

\bibitem{Lurie2009}
J.~Lurie,
\emph{On the classification of topological field theories},
in \emph{Current Developments in Mathematics 2008}, Int.\ Press, 2009, 129--280;
arXiv:0905.0465.

\bibitem{Hatcher2002}
A.~Hatcher,
\emph{Algebraic Topology},
Cambridge University Press, 2002.

\bibitem{MilnorStasheff}
J.~Milnor and J.~D. Stasheff,
\emph{Characteristic Classes},
Annals of Mathematics Studies \textbf{76}, Princeton University Press, 1974.

\bibitem{DijkgraafWitten}
R.~Dijkgraaf and E.~Witten,
\emph{Topological gauge theories and group cohomology},
Comm.\ Math.\ Phys.\ \textbf{129} (1990), 393--429.

\bibitem{Freed2014}
D.~S. Freed,
\emph{Anomalies and invertible field theories},
Proc.\ Symp.\ Pure Math.\ \textbf{88} (2015), 25--45; arXiv:1404.7224.

\bibitem{FreedHopkins}
D.~S. Freed and M.~J. Hopkins,
\emph{Reflection positivity and invertible topological phases},
Geom.\ Topol.\ \textbf{25} (2021), 1165--1330; arXiv:1604.06527.

\bibitem{BaezDolan1995}
J.~C. Baez and J.~Dolan,
\emph{Higher-dimensional algebra and topological quantum field theory},
J.\ Math.\ Phys.\ \textbf{36} (1995), 6073--6105.

\bibitem{BaezStay2010}
J.~C. Baez and M.~Stay,
\emph{Physics, topology, logic and computation: a Rosetta Stone},
in \emph{New Structures for Physics}, Lecture Notes in Physics \textbf{813}
(2010), 95--172; arXiv:0903.0340.

\bibitem{AtiyahSinger}
M.~F. Atiyah and I.~M. Singer,
\emph{The index of elliptic operators I},
Annals of Mathematics \textbf{87} (1968), 484--530.

\bibitem{Kitaev1997}
A.~Yu. Kitaev,
\emph{Fault-tolerant quantum computation by anyons},
Annals of Physics \textbf{303} (2003), 2--30; arXiv:quant-ph/9707021.

\bibitem{GradyPavlov2021}
D.~Grady and D.~Pavlov,
\emph{The geometric cobordism hypothesis},
arXiv:2111.01095 (2021).

\bibitem{KontsevichZagier2001}
M.~Kontsevich and D.~Zagier,
\emph{Periods},
in \emph{Mathematics Unlimited --- 2001 and Beyond}, Springer, 2001, 771--808.

\end{thebibliography}

\end{document}
